\begin{proof}[Proof of \cref{obj:prop:strict-enlargement}]
Fix an admissible strip exponent \(a\), and use the construction and selected propensity in \cref{obj:def:thin-strip} throughout.

\begin{enumerate}
\item Let the completed-data law and its treated response be
\[
P_{\mathrm{comp}}
:=
\mathcal L(X,A,Y(0),Y(1),Y),
\qquad
\mu_1(x):=Bp_0,\quad x\in\mathcal X.
\]
The construction gives
\[
0<p_0=\min\left\{\frac14,\frac{L}{4B}\right\}\le\frac14,
\qquad
Bp_0\le B\frac{L}{4B}=\frac L4.
\]
The radial term and strip increment are nonnegative, and truncation at one therefore gives
\[
0\le e_\star(x)\le1,\qquad x\in\mathcal X.
\]
Thus all conditional Bernoulli distributions are probability distributions. Integrating their product against the uniform covariate distribution shows that \(P_{\mathrm{comp}}\) and its observed-data marginal \(P_{\mathrm{strip}}\) are probability laws, with
\[
(P_{\mathrm{comp}})_X=(P_{\mathrm{strip}})_X
=\operatorname{Leb}_d\!\restriction_{\mathcal X}.
\]

The conditional distribution of \(Y(1)\) is the constant scaled Bernoulli distribution, so
\[
\mathbb E_{P_{\mathrm{comp}}}[Y(1)\mid X]
=(1-p_0)\,0+p_0B
=Bp_0
=\mu_1(X)
\quad\text{almost surely}.
\]
For the derivative order \(m\) in \cref{obj:def:holder}, the constant function satisfies
\[
\max_{|\alpha|\le m}\|D^\alpha\mu_1\|_\infty
=Bp_0\le\frac L4\le L,
\qquad
\max_{|\alpha|=m}\sup_{x\ne x'}
\frac{|D^\alpha\mu_1(x)-D^\alpha\mu_1(x')|}
{\|x-x'\|_\infty^{\beta-m}}
=0.
\]
Its derivatives have continuous boundary traces. Hence
\[
\mu_1\in\mathcal H^\beta(L),
\]
as defined in \cref{obj:def:holder}, and it supplies the response required by \cref{obj:ass:holder-response}.
% lean: CausalSmith.Stat.WeakOverlap.completionBind_holderResponse_constant_equal_arms

\item We verify the remaining restrictions in \cref{obj:def:model}. The uniform marginal gives
\[
P_{\mathrm{strip}}(X\in\mathcal X)=1,
\qquad
f(x)=1\ge c_f
\quad\text{for Lebesgue-almost every }x\in\mathcal X,
\]
which supplies \cref{obj:ass:density}. The completed construction gives
\[
Y=AY(1)+(1-A)Y(0)
\quad\text{almost surely},
\]
and, for almost every \(x\in\mathcal X\),
\[
\mathcal L\bigl((Y(0),Y(1)),A\mid X=x\bigr)
=
\mathcal L\bigl(Y(0),Y(1)\mid X=x\bigr)
\otimes\operatorname{Bernoulli}(e_\star(x)).
\]
These identities verify \cref{obj:ass:consistency,obj:ass:exchangeability}, and the treatment marginal gives
\[
P_{\mathrm{comp}}(A=1\mid X=x)=e_\star(x)
\quad\text{for almost every }x.
\]

For the overlap calculation, define the radial propensity and radial exponent by
\[
q_{\mathrm{rad}}(x)
:=F_R(R(x))^{1/(\gamma-1)},\quad x\in\mathcal X,
\qquad
\delta_{\mathrm{strip}}:=\frac d{\gamma-1}>0.
\]
The centered cube of radius \(r\) has side length \(2r\), whence
\[
F_R(r)=(2r)^d,\qquad 0\le r\le\frac12.
\]
In particular \(0\le q_{\mathrm{rad}}(x)\le1\), so the nonnegative strip increment gives
\[
q_{\mathrm{rad}}(x)\le e_\star(x),
\qquad
\{x\in\mathcal X:e_\star(x)\le t\}
\subseteq
\{x\in\mathcal X:F_R(R(x))\le t^{\gamma-1}\}.
\]

To bound the last event over its whole range, for \(0\le\tau_{\mathrm{cdf}}<1\) define
\[
r_{\mathrm{cdf}}(\tau_{\mathrm{cdf}})
:=\frac{\tau_{\mathrm{cdf}}^{1/d}}2.
\]
Then \(0\le r_{\mathrm{cdf}}(\tau_{\mathrm{cdf}})\le1/2\), and monotonicity of the \(d\)-th power on nonnegative reals gives
\[
\{x\in\mathcal X:F_R(R(x))\le\tau_{\mathrm{cdf}}\}
\subseteq
\{x\in\mathcal X:R(x)\le r_{\mathrm{cdf}}(\tau_{\mathrm{cdf}})\}.
\]
Consequently,
\[
P_{\mathrm{strip}}\{F_R(R(X))\le\tau_{\mathrm{cdf}}\}
\le
\bigl(2r_{\mathrm{cdf}}(\tau_{\mathrm{cdf}})\bigr)^d
=\tau_{\mathrm{cdf}}.
\]
This includes \(\tau_{\mathrm{cdf}}=0\); at \(\tau_{\mathrm{cdf}}=1\), the same upper bound follows from total probability one. Taking \(\tau_{\mathrm{cdf}}=t^{\gamma-1}\) therefore yields
\[
P_{\mathrm{strip}}\{e_\star(X)\le t\}
\le t^{\gamma-1}
\le Ct^{\gamma-1},
\qquad 0\le t\le1.
\]
This verifies \cref{obj:ass:global-tail}. At \(t=0\), the sharp bound and nonnegativity of \(e_\star\) also give
\[
P_{\mathrm{strip}}\{e_\star(X)=0\}=0,
\qquad
e_\star(X)>0\quad\text{almost surely}.
\]

The equal-arm product construction identifies the treated conditional distribution wherever treatment has positive conditional probability:
\[
\mathcal L_{P_{\mathrm{strip}}}(Y\mid A=1,X=x)
=\mathcal L\bigl(B\operatorname{Bernoulli}(p_0)\bigr),
\qquad
\mathbb E_{P_{\mathrm{strip}}}[Y\mid A=1,X=x]=Bp_0
\]
for \(P_X\)-almost every \(x\). Its support is contained in \([0,B]\). The bounded-range exponential-moment inequality, followed by enlargement of its exponent, gives
\[
\mathbb E_{P_{\mathrm{strip}}}
\!\left[\exp\{\lambda(Y-Bp_0)\}\mid A=1,X=x\right]
\le
\exp\left\{\frac{B^2\lambda^2}{8}\right\}
\le
\exp\left\{\frac{B^2\lambda^2}{2}\right\},
\qquad \lambda\in\mathbb R.
\]
Also \(|Bp_0|\le B\). For every \(\lambda\ne0\), it follows that
\[
\frac{
\sqrt{2\log\mathbb E_{P_{\mathrm{strip}}}
[\exp\{\lambda(Y-Bp_0)\}\mid A=1,X=x]}
}{|\lambda|}
\le B.
\]
The centered exponential moment is at least one by Jensen's inequality, so the square root is well defined. These estimates supply \cref{obj:ass:bounded-potential-outcomes}. Together with the smoothness already established and the stated parameter conditions, they prove
\[
P_{\mathrm{strip}}\in\mathcal P_\gamma.
\]
The uniform covariate marginal also supplies the asserted sharp bound
\[
P_{\mathrm{strip}}\{e_\star(X)\le t\}\le t^{\gamma-1},
\qquad t\in[0,1].
\]
% lean: CausalSmith.Stat.WeakOverlap.thinStrip_globalTailModel_of_holder
% lean: CausalSmith.Stat.WeakOverlap.stripPropensity_tail_of_uniform_covariateLaw

\item Fix arbitrary \(\rho,\nu,h_0>0\). For \(h>0\), define the centered Euclidean ball and local propensity supremum by
\[
\mathcal B_h:=\{x\in\mathbb R^d:\|x-x^\circ\|_2\le h\},
\qquad
M_{\mathrm{loc}}(h)
:=\sup_{x\in\mathcal B_h\cap\mathcal X}e_\star(x).
\]
For \(0<h\le1/2\), every coordinate displacement in \(\mathcal B_h\) has magnitude at most \(h\). On \(S\cap\mathcal B_h\), the strip inequality further gives
\[
|x_2-1/2|\le |x_1-1/2|^a\le h^a.
\]
Thus this intersection is contained in a rectangular box with second-coordinate side length \(2h^a\) and all other side lengths \(2h\). Uniformity implies
\[
P_{\mathrm{strip}}(X\in S\cap\mathcal B_h)
\le(2h^a)(2h)^{d-1}.
\]
For the denominator, the centered box with every coordinate displacement at most \(h/d\) lies in \(\mathcal X\) and in \(\mathcal B_h\), since
\[
\sum_{i=1}^d(x_i-1/2)^2
\le d(h/d)^2
=\frac{h^2}{d}\le h^2.
\]
Its volume therefore supplies
\[
P_{\mathrm{strip}}(X\in\mathcal B_h)
\ge(2h/d)^d>0.
\]
Dividing these bounds yields
\[
\frac{P_{\mathrm{strip}}(X\in S\cap\mathcal B_h)}
{P_{\mathrm{strip}}(X\in\mathcal B_h)}
\le
\frac{(2h^a)(2h)^{d-1}}{(2h/d)^d}
=d^d h^{a-1}.
\]
Because \(a-1>0\) and \(\delta_{\mathrm{strip}}>0\),
\[
d^d h^{a-1}\longrightarrow0,
\qquad
(2h)^{\delta_{\mathrm{strip}}}\longrightarrow0
\quad\text{as }h\downarrow0.
\]
Choose \(h\) satisfying
\[
0<h<\min\{h_0,1/2\},
\qquad
d^d h^{a-1}<\nu,
\qquad
(2h)^{\delta_{\mathrm{strip}}}<\rho/4.
\]

The center belongs to \(S\), and its radial term is zero. Hence
\[
e_\star(x^\circ)=\frac14,
\qquad
M_{\mathrm{loc}}(h)\ge\frac14.
\]
For \(x\in\mathcal B_h\setminus S\), the propensity equals its radial component and \(R(x)\le h\), giving
\[
e_\star(x)=q_{\mathrm{rad}}(x)
\le(2h)^{\delta_{\mathrm{strip}}}
<\rho/4
\le\rho M_{\mathrm{loc}}(h).
\]
Therefore
\[
\{x\in\mathcal B_h:e_\star(x)\ge\rho M_{\mathrm{loc}}(h)\}
\subseteq S\cap\mathcal B_h,
\]
and
\[
\frac{
P_{\mathrm{strip}}\{e_\star(X)\ge\rho M_{\mathrm{loc}}(h),
\ X\in\mathcal B_h\}
}{
P_{\mathrm{strip}}(X\in\mathcal B_h)
}
\le d^d h^{a-1}<\nu.
\]
This contradicts the numerical clause for the proposed triple. Since the triple was arbitrary, the numerical anti-concentration clause fails. The argument uses the pointwise bound \(e_\star\le1\), so it applies with propensity values equal to one admitted.
% lean: CausalSmith.Stat.WeakOverlap.thinStripLaw_not_dornA3AntiConcentration

\item The full singleton-family condition in \cref{obj:def:dorn-a3} entails its numerical clause with the same constants \(\rho,\nu,h_0\). Thus the preceding contradiction gives
\[
\neg\mathsf A_3^{\mathrm{Dorn}}(P_{\mathrm{strip}})
\]
for the selected propensity \(e_\star\).
% lean: CausalSmith.Stat.WeakOverlap.dornA3AntiConcentrationOnCube_of_dornA3OnCube

\item For the design limit, define, for every \(h>0\),
\[
\mathcal Q_{\mathrm{ctr},h}
:=\{x\in\mathbb R^d:R(x)\le h\},
\]
\[
\begin{aligned}
D_{\mathrm{strip}}(h)
&:=P_{\mathrm{strip}}\{A=1,\ X\in\mathcal Q_{\mathrm{ctr},h}\},\\
N_{\mathrm{strip}}(h)
&:=\int_{\{A=1,\ X\in\mathcal Q_{\mathrm{ctr},h}\}}
\left(\frac{X_2-1/2}{h}\right)^2\,dP_{\mathrm{strip}},\\
T_{\mathrm{strip}}(h)
&:=\frac{N_{\mathrm{strip}}(h)}{D_{\mathrm{strip}}(h)}.
\end{aligned}
\]
For this definition, a quotient with zero denominator is assigned value zero.

Introduce the positive constants and envelope
\[
b_{\mathrm{strip}}
:=\frac{1}{2(4d)^{a+d-1}},
\qquad
c_{\mathrm{rad}}:=2^{d+\delta_{\mathrm{strip}}},
\qquad
c_{\mathrm{strip}}:=2^d,
\]
\[
G_{\mathrm{strip}}(h)
:=
\frac{
c_{\mathrm{rad}}h^{d+\delta_{\mathrm{strip}}}
+c_{\mathrm{strip}}h^{d+3a-3}
}{
b_{\mathrm{strip}}h^{d+a-1}
},
\qquad h>0.
\]
Factoring powers of \(h\) gives
\[
G_{\mathrm{strip}}(h)
=
\frac{c_{\mathrm{rad}}}{b_{\mathrm{strip}}}
h^{1+\delta_{\mathrm{strip}}-a}
+
\frac{c_{\mathrm{strip}}}{b_{\mathrm{strip}}}
h^{2a-2}.
\]
The strip-exponent restrictions give
\[
1+\delta_{\mathrm{strip}}-a>0,
\qquad
2a-2>0,
\]
so both powers tend to zero and
\[
G_{\mathrm{strip}}(h)\longrightarrow0
\quad\text{as }h\downarrow0.
\]
% lean: CausalSmith.Stat.WeakOverlap.thinStrip_raw_second_envelope_tendsto_zero

\item We bound the moment by this envelope. Suppose \(0<h\le1/2\). On \(\mathcal Q_{\mathrm{ctr},h}\), the normalized second-coordinate square is at most one. If \(x\in S\cap\mathcal Q_{\mathrm{ctr},h}\), the strip inequality and \(e_\star(x)\le1\) give
\[
\left(\frac{x_2-1/2}{h}\right)^2e_\star(x)
\le\left(\frac{h^a}{h}\right)^2.
\]
If \(x\in\mathcal Q_{\mathrm{ctr},h}\setminus S\), the radial bound gives
\[
\left(\frac{x_2-1/2}{h}\right)^2e_\star(x)
\le e_\star(x)\le(2h)^{\delta_{\mathrm{strip}}}.
\]
Combining these two cases and integrating,
\[
\begin{aligned}
&\int_{\mathcal Q_{\mathrm{ctr},h}}
\left(\frac{x_2-1/2}{h}\right)^2e_\star(x)\,dx\\
&\qquad\le
(2h)^{\delta_{\mathrm{strip}}}
P_{\mathrm{strip}}(X\in\mathcal Q_{\mathrm{ctr},h})
+
\left(\frac{h^a}{h}\right)^2
P_{\mathrm{strip}}(X\in S\cap\mathcal Q_{\mathrm{ctr},h}).
\end{aligned}
\]
Here \(\mathcal Q_{\mathrm{ctr},h}\subseteq\mathcal X\), and its volume is \((2h)^d\). The same rectangular enclosure used above bounds the strip intersection, so
\[
\begin{aligned}
&\int_{\mathcal Q_{\mathrm{ctr},h}}
\left(\frac{x_2-1/2}{h}\right)^2e_\star(x)\,dx\\
&\qquad\le
(2h)^{\delta_{\mathrm{strip}}}(2h)^d
+\left(\frac{h^a}{h}\right)^2(2h^a)(2h)^{d-1}\\
&\qquad=
c_{\mathrm{rad}}h^{d+\delta_{\mathrm{strip}}}
+c_{\mathrm{strip}}h^{d+3a-3}.
\end{aligned}
\]

For the treated-mass lower bound, define
\[
r_{\mathrm{core}}(h):=\frac{h}{4d},
\]
\[
\mathcal K_h
:=
\left\{
x\in\mathbb R^d:
|x_2-1/2|\le r_{\mathrm{core}}(h)^a,\
r_{\mathrm{core}}(h)\le x_i-1/2\le2r_{\mathrm{core}}(h)
\ \text{for every }i\ne2
\right\}.
\]
Since \(0<r_{\mathrm{core}}(h)\le1/8\) and \(a>1\),
\[
r_{\mathrm{core}}(h)^a\le r_{\mathrm{core}}(h).
\]
The coordinate bounds place \(\mathcal K_h\) in \(\mathcal X\), and they imply
\[
|x_2-1/2|
\le r_{\mathrm{core}}(h)^a
\le |x_1-1/2|^a,
\]
\[
\|x-x^\circ\|_2^2
\le d\bigl(2r_{\mathrm{core}}(h)\bigr)^2
=\frac{h^2}{4d}\le h^2.
\]
Thus
\[
\mathcal K_h\subseteq S\cap\mathcal B_h
\subseteq S\cap\mathcal Q_{\mathrm{ctr},h}.
\]
Its second-coordinate side length is \(2r_{\mathrm{core}}(h)^a\), and each other side length is \(r_{\mathrm{core}}(h)\). Since \(e_\star\ge1/4\) on \(S\),
\[
\begin{aligned}
\int_{\mathcal Q_{\mathrm{ctr},h}}e_\star(x)\,dx
&\ge\int_{\mathcal B_h\cap\mathcal X}e_\star(x)\,dx\\
&\ge\frac14
\bigl(2r_{\mathrm{core}}(h)^a\bigr)
r_{\mathrm{core}}(h)^{d-1}\\
&=b_{\mathrm{strip}}h^{d+a-1}>0.
\end{aligned}
\]

Finally, conditioning on \(X\) and using
\(P_{\mathrm{strip}}(A=1\mid X)=e_\star(X)\) gives the numerator and denominator identities
\[
N_{\mathrm{strip}}(h)
=
\int_{\mathcal Q_{\mathrm{ctr},h}}
\left(\frac{x_2-1/2}{h}\right)^2e_\star(x)\,dx,
\qquad
D_{\mathrm{strip}}(h)
=
\int_{\mathcal Q_{\mathrm{ctr},h}}e_\star(x)\,dx.
\]
All integrands are nonnegative. Dividing the numerator upper bound by the positive denominator lower bound therefore yields
\[
0\le T_{\mathrm{strip}}(h)\le G_{\mathrm{strip}}(h),
\qquad 0<h\le1/2.
\]
The squeeze theorem proves
\[
\lim_{h\downarrow0}T_{\mathrm{strip}}(h)=0,
\]
which is exactly the asserted normalized treated-design limit.
% lean: CausalSmith.Stat.WeakOverlap.stripSecondMoment_tendsto_zero

\item The law \(P_{\mathrm{strip}}\) itself supplies the final existence claim: it belongs to \(\mathcal P_\gamma\), and its full singleton-family Dorn predicate fails for \(e_\star\).
% lean: CausalSmith.Stat.WeakOverlap.thinStrip_mem_model_not_dornA3
\end{enumerate}
\end{proof}
